% ============================================================================
%  04-viet-va-viet-lai — Viết như công cụ nghĩ, và viết lại ở đúng tầng
%  Ngày tạo: 2026-09-09 | Sửa gần nhất: 2026-09-09 | Ôn gần nhất: chưa
%  Dependency: 02-moi-truong-phan-hoi. Mức độ: nên nắm (phần Boice/Sword: bằng chứng yếu).
% ============================================================================
\documentclass[../hieu-suat.tex]{subfiles}
\begin{document}
\chapter{Viết lại: người mới sửa chữ, người thành thục sửa lập luận}
\ptrtarget{04}\ptrtarget{04-viet-va-viet-lai}
\dependency{\ptr{02} --- viết là cách lấy phản hồi từ chính suy nghĩ của mình khi chưa có ai để hỏi. Đọc gần \ptr{05}.}

\begin{tomtat}
Có một khác biệt đã được quan sát trực tiếp giữa người viết mới và người viết thành thục, và nó không phải vốn từ: họ \emph{viết lại ở hai tầng khác nhau}. Người mới sửa chữ và sửa lỗi; người thành thục vứt bỏ và dựng lại cấu trúc lập luận, coi bản nháp đầu là suy nghĩ thô chứ không phải sản phẩm gần xong. Chương này lấy kết quả đó làm trục, ghép thêm mô hình quá trình viết, quy tắc cấu trúc bài báo khoa học, và một tranh cãi thật về việc ``viết mỗi ngày'' --- nơi tôi sẽ chỉ ra rằng bằng chứng cho lời khuyên phổ biến nhất trong giới học thuật là rất yếu. Nếu chỉ nhớ một điều: đừng đánh bóng câu chữ trước khi lập luận đã ổn định, vì mỗi câu được đánh bóng làm tăng cái giá tâm lý của việc xoá nó.
\end{tomtat}

\begin{nentang}
\kn{bản nháp không (zero draft)}{bản viết ra chỉ để nghĩ, được chấp nhận là xấu, và không ai ngoài tác giả đọc.}
\kn{vị trí chủ đề và vị trí nhấn (topic position, stress position)}{đầu câu và cuối câu --- hai chỗ mà người đọc theo mặc định đi tìm ``câu này nói về cái gì'' và ``điều quan trọng là gì'' \cite{gopen1990}.}
\end{nentang}

\begin{khoigoi}
\h{Viết lại nghĩa là gì --- sửa câu cho mượt, hay việc khác hẳn?}
\h{Vì sao viết ra lại làm lộ chỗ mình chưa hiểu, trong khi nghĩ trong đầu thì không?}
\h{``Viết mỗi ngày'' có phải lời khuyên có bằng chứng không?}
\end{khoigoi}

\section{Kết quả gốc: hai tầng viết lại khác nhau}

Sommers \cite{sommers1980} quan sát tám sinh viên năm nhất và bảy người viết chuyên nghiệp cùng làm một việc: viết ba bài, mỗi bài viết lại hai lần, kèm phỏng vấn. Điều bà tìm thấy không phải là người chuyên nghiệp viết lại \emph{nhiều hơn}, mà là họ viết lại \emph{ở tầng khác} \nD{}. Từ vựng mà sinh viên dùng để mô tả việc viết lại xoay quanh việc thay chữ và sửa lỗi --- ``đổi từ'', ``sửa lỗi'' --- và họ đối xử với bản nháp đầu như một sản phẩm gần hoàn chỉnh chỉ cần đánh bóng. Người viết chuyên nghiệp mô tả cùng hoạt động đó bằng những từ khác hẳn: xem xét lại, nghĩ lại, khám phá; họ chú ý tới sự phát triển của lập luận, tới tổ chức, và tới việc người đọc sẽ hiểu ra sao; họ sẵn sàng cắt cả đoạn lớn và sinh ý mới giữa chừng.

\begin{trucgiac}
Cơ chế đằng sau chênh lệch này, theo cách tôi hiểu \nM{}: mỗi câu đã được đánh bóng là một khoản đầu tư, và khoản đầu tư đã bỏ ra làm việc xoá nó đắt hơn về mặt tâm lý. Người sửa câu trước rồi mới xét cấu trúc sau đã tự khoá mình vào cấu trúc hiện có --- không phải vì nó đúng, mà vì phá nó bây giờ tốn quá nhiều thứ vừa làm xong. Người viết thành thục tránh cái bẫy đó bằng cách giữ bản nháp đầu ở trạng thái rẻ tiền, cố ý xấu, dễ vứt.
\end{trucgiac}

\section{Viết là quá trình đệ quy, không phải dây chuyền}

Flower và Hayes \cite{flower1981} mô hình hoá việc viết thành ba quá trình --- lập kế hoạch, chuyển thành chữ, và xem lại --- cộng một cơ chế \emph{giám sát} quyết định khi nào chuyển giữa chúng. Điểm quan trọng của mô hình là ba quá trình này lồng vào nhau và gọi nhau ở bất kỳ thời điểm nào, chứ không xếp thành ba giai đoạn nối tiếp \nT{}. Điều này giải thích vì sao ``viết xong rồi mới sửa'' là một mô tả sai về cách viết thật, và vì sao việc \emph{viết} lại làm thay đổi \emph{điều mình nghĩ}: khâu chuyển thành chữ liên tục gửi phản hồi ngược về khâu lập kế hoạch.

Đó cũng là lý do viết là công cụ nghiên cứu chứ không phải khâu cuối \nM{}. Một lập luận chưa viết ra tồn tại trong đầu dưới dạng cảm giác mạch lạc; khi bị buộc phải thành câu có chủ ngữ và động từ, những chỗ đứt gãy hiện ra ngay --- đó chính là loại phản hồi nhanh mà chương \ptr{02} nói là môi trường nghiên cứu không tự cấp.

\section{Cấu trúc: một bài báo mang đúng một thông điệp}

Whitesides \cite{whitesides2004} mô tả một quy trình đơn giản và hiếm ai làm: bắt đầu bằng \emph{dàn ý} ngay từ lúc khởi động dự án, chứ không phải khi đã có kết quả; dàn ý đó gồm các hình và bảng dự kiến, và nó chính là bản kế hoạch nghiên cứu \nT{}. Hệ quả rất thực dụng: khi dàn ý không dựng nổi thành một lập luận, đó là dấu hiệu thí nghiệm còn thiếu, và biết điều đó trước khi chạy thì rẻ hơn nhiều so với biết sau.

Mensh và Kording \cite{mensh2017} bổ sung phần cấu trúc ở ba tỉ lệ: toàn bài, đoạn, và trong một đoạn; mỗi tỉ lệ đều theo mạch bối cảnh--nội dung--kết luận, và phần tóm tắt là chỗ duy nhất chứa cả ba \nT{}. Quy tắc mà tôi thấy đắt giá nhất trong bài đó: một bài báo mang \emph{một} thông điệp. Khi có hai, người đọc không nhớ cái nào.

Ở tầng câu, Gopen và Swan \cite{gopen1990} lập luận rằng vấn đề của văn khoa học khó đọc thường không phải nội dung phức tạp mà là thông tin bị đặt sai chỗ so với nơi người đọc chờ đợi nó: người đọc lấy \emph{chủ đề} của câu ở đầu câu và tìm \emph{thông tin quan trọng nhất} ở cuối câu \nT{}. Đây là công cụ cho lượt viết lại cuối, và chỉ cho lượt cuối.

\section{Một tranh cãi thật: ``viết mỗi ngày''}

Lời khuyên phổ biến nhất trong giới học thuật là viết mỗi ngày trong những phiên ngắn. Nguồn gốc của nó là các nghiên cứu của Boice. Trong nghiên cứu can thiệp một năm \cite{boice1989}, ba nhóm giảng viên trẻ: nhóm một giữ nguyên thói quen cũ (dồn việc viết vào những khối thời gian lớn hiếm khi xuất hiện), nhóm hai viết mỗi ngày khoảng 30 phút và ghi nhật ký thời gian, nhóm ba làm như nhóm hai cộng thêm việc bị kiểm tra đột xuất hai lần mỗi tuần. Cuối năm, nhóm ba sản xuất trung bình hơn gấp đôi nhóm hai và hơn chín lần nhóm đối chứng \nD{} --- theo tường thuật chi tiết của Sword \cite[tr.~315]{sword2016}. Trong một thí nghiệm khác \cite{boice1983}, 27 giảng viên chia ba nhóm chín người --- kiêng viết, viết khi có hứng, và viết theo lịch có chế tài --- cho sản lượng trung bình lần lượt khoảng 0,2 / 0,8 / 3,2 trang mỗi ngày \nD{}.

Sword \cite{sword2016} không bác bỏ các con số đó, bà kiểm tra nền của chúng, và đây là chỗ đáng học về cách phản biện một bằng chứng \nM{}. Thứ nhất, cỡ mẫu rất nhỏ (chín người mỗi nhóm), toàn bộ đến từ vài trường ở cùng một địa phương, và toàn là người tự nhận mình hay trì hoãn. Thứ hai, điều kiện của nhóm đối chứng rất nhân tạo: họ bị yêu cầu \emph{kiêng} mọi việc viết không thiết yếu trong mười tuần, nên con số 0,2 trang không đo được gì về người bình thường. Thứ ba, 15 trong số 42 người tham gia ban đầu đã bỏ cuộc, phần lớn vì thấy việc ghi chép phiền --- và tỉ lệ bỏ cuộc này không được nhắc tới khi chính tác giả tóm tắt lại nghiên cứu trong sách của mình \nD{}.

Rồi bà đưa dữ liệu ngược chiều: phỏng vấn 100 nhà nghiên cứu thành công ở 45 trường đại học trên bốn châu lục trong bốn năm, chỉ \textbf{12,8\%} nhóm được đồng nghiệp giới thiệu là ``người viết xuất sắc'' và \textbf{11,5\%} nhóm dự các khoá phát triển kỹ năng viết nói rằng họ viết mỗi ngày; chỉ hai trong số 100 người mô tả một nếp viết gần giống mô hình được khuyến nghị \nD{}. Bà cũng lưu ý một con số từ chính Boice: ước lượng chỉ khoảng 12\% người làm nghiên cứu thực sự bị chứng tắc viết \nD{}, trong khi gần như toàn bộ chương trình can thiệp được thiết kế cho nhóm đó.

\begin{cambay}
Đây là ví dụ mẫu về cách một lời khuyên trở thành hiển nhiên mà không cần bằng chứng mạnh: một nghiên cứu nhỏ, điều kiện nhân tạo, được tóm tắt lại nhiều lần trong sách phổ thông, mỗi lần mất thêm một phần các giới hạn, cho tới khi nó thành quy tắc đạo đức nghề nghiệp. Khi gặp một lời khuyên ai cũng nhắc lại, việc đáng làm là truy về bài gốc và xem cỡ mẫu \nM{}.
\end{cambay}

Phần còn lại sau khi trừ đi những gì không đứng vững \nM{}: bằng chứng \emph{không} ủng hộ mệnh lệnh ``phải viết mỗi ngày''. Thứ có vẻ đứng vững hơn là ba điều yếu hơn nhưng thực dụng --- đừng chờ khối thời gian lớn và cảm hứng, vì chúng hiếm khi đến; ghi lại thời gian viết thật để có dữ liệu về chính mình thay vì cảm giác; và một dạng trách nhiệm với người khác giúp duy trì nếp làm việc (nhóm có chế tài luôn cao hơn nhóm không có trong cả hai thí nghiệm của Boice). Nhịp cụ thể --- ngày nào, dài bao lâu --- là thứ tôi phải tự đo trên chính mình, và chương \ptr{10} nói cách đo.

\section{Quy trình viết lại nhiều tầng}

Ghép Sommers, Flower--Hayes, Whitesides, Mensh--Kording và Gopen--Swan lại, tôi dùng sáu lượt, \textbf{đi từ tầng cao xuống tầng thấp, không bao giờ ngược lại} \nM{}. Quy tắc chặn: chỉ được xuống tầng dưới khi tầng trên đã ổn định, vì mọi công sức bỏ vào tầng dưới đều mất khi tầng trên đổi.

\begin{enumerate}
\item \textbf{Lượt 0 --- dàn ý của các khẳng định.} Một trang: thông điệp duy nhất của bài, các khẳng định phụ, và \emph{bằng chứng nào} chống lưng cho từng khẳng định. Làm ngay khi bắt đầu dự án \cite{whitesides2004}, không phải khi có kết quả.
\item \textbf{Lượt 1 --- bảng bằng chứng.} Với mỗi khẳng định: thí nghiệm nào, số nào, điều kiện nào, có thiếu nhóm đối chứng không. Chỗ trống trong bảng này là danh sách thí nghiệm còn phải chạy.
\item \textbf{Lượt 2 --- cấu trúc.} Sắp lại thứ tự mục sao cho mỗi mục trả lời một câu hỏi mà mục trước vừa làm nảy sinh. Ở lượt này việc cắt bỏ cả mục là bình thường, và đó là lý do chưa được đánh bóng câu.
\item \textbf{Lượt 3 --- đoạn.} Mỗi đoạn một ý; câu đầu đoạn nói ý đó. Đọc riêng chuỗi câu đầu đoạn: nếu chuỗi đó không thành một lập luận liền mạch thì bài chưa xong \nM{}.
\item \textbf{Lượt 4 --- câu.} Bây giờ mới tới vị trí chủ đề và vị trí nhấn \cite{gopen1990}: chủ đề của câu nằm đầu, thông tin mới và quan trọng nằm cuối.
\item \textbf{Lượt 5 --- đối kháng.} Đọc bài như một người phản biện muốn bác bỏ nó; đây là nội dung của chương \ptr{05}.
\end{enumerate}

\begin{cannho}
Viết lại không phải là sửa chữ. Nếu sau một buổi viết lại mà cấu trúc lập luận không đổi và không có đoạn nào bị xoá, khả năng cao là tôi vừa đánh bóng chứ chưa viết lại \nM{}.
\end{cannho}

\begin{traloi}
\tl{Ở người thành thục, viết lại là dựng lại lập luận và tổ chức, thậm chí cắt bỏ và sinh ý mới; ở người mới, nó là thay chữ và sửa lỗi \cite{sommers1980}.}
\tl{Vì khâu chuyển ý thành chữ liên tục gửi phản hồi ngược về khâu lập kế hoạch \cite{flower1981}; cảm giác mạch lạc trong đầu không chịu được ràng buộc của một câu hoàn chỉnh.}
\tl{Không, theo nghĩa mệnh lệnh. Bằng chứng gốc là các nghiên cứu cỡ mẫu chín người với nhóm đối chứng nhân tạo và tỉ lệ bỏ cuộc cao \cite{boice1983,boice1989,sword2016}, và khảo sát 100 nhà nghiên cứu thành công cho thấy chỉ khoảng 12\% viết mỗi ngày \cite{sword2016}.}
\end{traloi}

\begin{tukiem}
\item Nêu bốn tầng của việc viết lại theo thứ tự, và giải thích vì sao đi ngược thứ tự lại tốn kém.
\item Trong bản thảo gần nhất của bạn, chuỗi câu đầu đoạn có tự nó thành một lập luận không? Nếu không, chỗ nào đứt?
\item Bằng chứng cho ``viết mỗi ngày'' yếu ở ba điểm nào? Trả lời không nhìn bài.
\item Dàn ý của Whitesides được viết \emph{khi nào}, và vì sao thời điểm đó quan trọng hơn nội dung dàn ý?
\end{tukiem}

\begin{onbai}
\ar[3]{Sommers 1980}{Người mới sửa chữ và lỗi; người thành thục dựng lại lập luận, coi nháp đầu là suy nghĩ thô.}
\ar[3]{Thứ tự viết lại}{Khẳng định \(\rightarrow\) bằng chứng \(\rightarrow\) cấu trúc \(\rightarrow\) đoạn \(\rightarrow\) câu \(\rightarrow\) đối kháng. Không đi ngược.}
\ar[3]{Viết là để nghĩ}{Ba quá trình lồng nhau có giám sát \cite{flower1981}; viết ra làm lộ chỗ đứt của lập luận.}
\ar[2]{Dàn ý từ đầu dự án}{Dàn ý gồm hình và bảng dự kiến chính là kế hoạch nghiên cứu \cite{whitesides2004}.}
\ar[2]{Một bài một thông điệp}{Bối cảnh--nội dung--kết luận ở ba tỉ lệ \cite{mensh2017}.}
\ar[2]{Vị trí chủ đề / vị trí nhấn}{Chủ đề đầu câu, thông tin quan trọng cuối câu \cite{gopen1990}; chỉ dùng ở lượt cuối.}
\ar[2]{``Viết mỗi ngày''}{Bằng chứng yếu: n=9/nhóm, đối chứng nhân tạo, 15/42 bỏ cuộc; chỉ 12,8\% người thành công làm vậy \cite{sword2016}.}
\ar[1]{Cái còn lại sau phản biện}{Đừng đợi cảm hứng và khối lớn; ghi thời gian thật; có người cùng chịu trách nhiệm.}
\end{onbai}

\end{document}
